\label{sec:resumo_plano}

The project, which ran from 1 September 2025 to 31 August 2026, continued the previous stage. In that stage, a dense neural network of \num{8195} parameters estimated the interaction vertex of \qty{10}{\MeV} deposits in the third far detector module (FD3) of the DUNE experiment (Deep Underground Neutrino Experiment) from \num{60} descriptors of the scintillation light \citep{vicentin2024study}.

The objective was to use machine learning to estimate the three-dimensional position of energy deposits of the order of MeV from the light recorded by a photon detection system of the APEX type (Aluminum Profiles with Embedded X-ARAPUCA), as a measure of the performance of that system. The plan envisaged measuring the spatial resolution in each coordinate and comparing it with that of the analytic estimators of the APEX study \citep{marinho2025apex}. The reconstruction would also be extended to other deposited energies and to other sensor geometries, as the conclusion of the previous stage proposed.

The methodology started from the simulation used in that stage, with the photon counts rescaled from \qty{100}{\MeV} to \qty{10}{\MeV} and binomial detection with an efficiency of \qty{2}{\percent}. The plan included a hyperparameter search, an evaluation on events held out for testing in terms of the mean absolute error (MAE), the root-mean-square error (RMSE) and the coefficient of determination ($R^2$) along each axis, and a study as a function of energy and efficiency. The expected outcome was a per-coordinate resolution better than that of the previous stage, whose test MAE ranged from \num{12.03} to \qty{18.18}{\cm} per axis \citep{vicentin2024study}, and better than that of the analytic estimators, with a model lightweight enough for use in real-time data acquisition. The results would be presented at the 2026 Scientific Initiation Meeting of the Aeronautics Institute of Technology (ENCITA) and in this report.
